\section{Методы}

\subsection{Программные средства}

Для вычислений использовались \texttt{Python}, \texttt{NumPy},
\texttt{SciPy} и \texttt{Matplotlib}. \texttt{SciPy} применялся для
построения выпуклой оболочки алгоритмом \texttt{ConvexHull},
\texttt{NumPy}~--- для генерации точек и векторизованной проверки
принадлежности кругам, \texttt{Matplotlib}~--- для построения и сохранения
графиков.

\newpage

\subsection{Задача 1: выпуклая оболочка}

Генерируются 20 точек с координатами из диапазона $[-5,5]$. Функция
\texttt{ConvexHull} возвращает индексы вершин оболочки в порядке
обхода. Внутренние точки обозначены синим прозрачным цветом, а вершины
оболочки~--- красным. Сама оболочка закрашена полупрозрачным
зеленым цветом.

\inputmintedbox{python}{code/lab02-1.py}

\newpage

\subsection{Задача 2: пересечение кругов}

Рассматриваются множества
\[
	A=\{(x,y):x^2+y^2\leq 3^2\},
\]
\[
	B=\{(x,y):(x-2)^2+(y-1)^2\leq 2.5^2\}.
\]
На сетке $[-5,5]\times[-5,5]$ с шагом $0.01$ для каждой точки вычисляются логические маски принадлежности $A$, $B$ и $A\cap B$. Три непересекающиеся категории отображаются разными цветами.

\inputmintedbox{python}{code/lab02-2.py}

\subsection{Задача 3: критерий по ориентации}

Функция \texttt{is\_convex} последовательно рассматривает все тройки соседних вершин, включая переход через последнюю вершину. Для каждой тройки вычисляется знак ориентированного произведения. Первый ненулевой знак сохраняется как эталонный; если встречается противоположный знак, функция немедленно возвращает \texttt{False}. В тестах использованы выпуклый многоугольник с коллинеарной точкой, невыпуклый многоугольник с выемкой и звёздный многоугольник.

\inputmintedbox{python}{code/lab02-3.py}
\begin{mintedbox}{python}
# Одна точка лежит на линии между двумя другими точками,
# но не нарушает выпуклость.
# Повторяющиеся точки не нарушают выпуклость
oriented_points_convex = np.array(
	[[0, 0], [1, 0], [1, 1], [0.5, 1], [0, 1], [0, 0]]
	)
oriented_points_non_convex = np.array(
	[[0, 0], [1, 0], [0.5, 0.5], [1, 1], [0, 1], [0, 0]]
	)
star_points = np.array(...)
print(is_convex(oriented_points_convex))  # True
print(is_convex(oriented_points_non_convex))  # False
print(is_convex(star_points))  # False
\end{mintedbox}
